% !TeX root = ../main.tex

\xchapter{总结与展望}{Conclusion and Future Work}

\xsection{全文总结}{Summary}

本文围绕深度强化学习在真实部署场景中面临的观测分布偏移问题，分别针对高维视觉输入和低维结构化状态两类典型范式，提出了各具针对性的鲁棒表征学习方法。两项工作在问题动机上高度统一——均以"如何在训练与测试观测条件不一致时维持策略的泛化性能"为核心命题——并分别从数据增强策略优化和自监督缺失重建两条互补技术路径给出了解答。


视觉强化学习以原始像素图像作为观测输入，现有数据增强方法在提升跨视觉域泛化能力的同时，往往因强增强操作破坏状态序列的马尔可夫平滑性而引发 TD 自举方差急剧攀升，导致训练稳定性与泛化性能难以兼顾。PICK 的核心思路是将增强方法的选择从任务无关的随机过程转变为由 Critic 梯度信号端到端驱动的任务感知学习过程。具体而言，PICK 引入独立的 PickNetwork 模块，通过将高维视觉特征压缩为数值稳定的任务感知编码并经 Softmax 输出增强概率分布，使智能体能够根据当前任务的奖励结构动态抑制破坏性增强、优先选用与任务语义相适配的增强策略。与此同时，非对称 Actor-Critic 架构通过限制目标 $Q$ 值仅在原始未增强观测流上估计，从根本上切断了增强随机性对 Bellman 目标的方差污染。本文还引入随机频率掩盖与随机缩放两种新型增强方法，补充了频率扰动与尺度变换两类候选维度。实验结果表明，PICK 在 DMControl 六项连续控制任务上取得最优样本效率，并在 DMC-GB2 的高难度分布外测试环境中全面超越 DrQ、RAD、SVEA、SADA 等先进基线，验证了任务感知增强选择在主动塑造策略泛化能力方面的有效性。


离线强化学习在测试阶段面临的核心挑战之一是观测维度的部分缺失：训练时基于完整状态学习的策略在测试时遭遇随机维度丢失或持续维度缺失时，性能往往急剧退化。Masked-IQL 针对这一问题构建了从"残缺观测"到"鲁棒决策"的完整信息处理路径。可学习掩码标记机制通过引入独立的可训练嵌入替代零值填充，从根本上消除了真实零值与缺失标记之间的语义歧义。随机掩码比例采样策略使编码器在训练阶段同时暴露于不同缺失比例，实现了对测试时未知缺失率的跨域泛化。残缺-完整表征一致性损失配合编码器 LayerNorm，在连续潜在空间中主动对抗高缺失率下的表征分布偏移。轻量历史窗口聚合模块以极低的计算开销聚合多步历史编码，使相邻时间步中被观测到的维度能够对当前缺失维度进行隐式互补。两阶段解耦训练框架将掩码重建训练限定于表征学习阶段，确保 IQL 决策模块在稳定的完整状态编码上收敛。在 D4RL 基准的三个 MuJoCo 环境上，针对动态缺失与因子消减两种场景的系统性实验表明，Masked-IQL 在多数设定下取得最优或极具竞争力的泛化性能，性能退化曲线相较于对比方法更为平滑稳定。消融实验进一步揭示了两个核心组件在不同场景下互补的重要性规律：可学习 Mask Token 在动态缺失场景中贡献约 14.9\% 的一致性性能增益，历史窗口聚合在因子消减场景中则成为防止性能崩塌的决定性因素。

综上，本文两项工作分别从任务感知数据增强选择与自监督缺失表征学习两个维度，为深度强化学习在观测分布偏移下的鲁棒泛化提供了具有互补性的方法论贡献。


\xsection{研究展望}{Future Work}

本文的两项工作均在各自的问题设定中取得了显著进展，但仍存在若干值得深入探索的开放性问题，以下从方法延伸与场景拓展两个层面进行展望。

\xsubsection{PICK 方法的后续拓展}{Extensions of PICK}

增强方法的连续参数自适应调节。当前 PickNetwork 的选择粒度为增强方法级别的离散概率分配，尚未实现对单一增强方法内部参数（如旋转角度范围、频率掩盖强度、卷积核方差）的连续自适应调节。未来可借鉴超网络（Hypernetwork）或参数化策略梯度的思路，将 PickNetwork 扩展为同时输出离散增强选择与连续增强强度的联合分布，进一步提升增强策略与任务状态的精细适配程度。

跨任务迁移与多任务共享选择策略。现有 PICK 为每个任务独立训练 PickNetwork，无法利用跨任务的增强有效性先验。在多任务强化学习场景下，可引入共享主干与任务特定适配层的结构，使 PickNetwork 在多任务联合训练中习得可迁移的增强偏好元知识，从而在新任务上实现更快的增强策略收敛。

与自监督表征学习的深度融合。PICK 当前独立于表征学习过程运行，未来可探索将增强选择信号与对比学习、掩码自编码等自监督目标联合优化，使增强方法的选择同时服务于策略鲁棒性与表征语义一致性两个目标，有望在更少的在线交互下实现更强的分布外泛化。

\xsubsection{Masked-IQL 方法的后续拓展}{Extensions of Masked-IQL}

更复杂缺失模式的泛化。本文主要考察了动态随机缺失与持续维度缺失两种典型场景，而真实传感器故障往往呈现出更复杂的时序相关缺失模式（如突发性批量缺失、周期性维度交替缺失）。未来可设计针对时序相关缺失的专用掩码采样策略，或引入缺失模式的潜变量建模，使编码器对任意缺失结构具备更广泛的泛化能力。

向在线与半离线强化学习的扩展。当前 Masked-IQL 的两阶段框架针对纯离线设定设计。若将其扩展至在线微调场景，可探索在少量在线交互中以自监督方式持续更新编码器，在不破坏离线预训练表征的前提下逐步适应测试环境的真实缺失分布，实现离线预训练与在线适应的有机结合。

与多模态观测及模型基础强化学习的结合。在具身智能等复杂场景中，智能体通常同时接收视觉图像、关节状态、触觉信号等多模态观测，任意模态的缺失均会影响决策质量。Masked-IQL 的掩码重建思路天然适合推广至多模态部分可观测设定，通过跨模态一致性损失对齐不同缺失组合下的表征，可为多传感器融合系统的鲁棒部署提供理论支撑。与此同时，将缺失鲁棒编码器嵌入基于模型的强化学习框架（如 MBPO、DreamerV3）的世界模型，有望在残缺观测下同样实现准确的环境动力学预测，为残缺传感器场景下的规划与想象提供基础能力。

\xsubsection{统一视角下的未来方向}{Unified Perspective on Future Directions}

从更宏观的视角来看，本文两项工作分别面向在线视觉强化学习与离线低维状态强化学习，但均指向同一个核心问题：如何让强化学习智能体在训练与部署之间的观测条件不一致时保持稳健。随着大规模预训练视觉-语言模型和强化学习基础模型的兴起，探索通用预训练表征在缺失鲁棒性方面的潜力，以及如何将任务感知的数据增强选择与自监督掩码重建机制融合进统一的基础模型微调框架，将是该研究方向未来最具前景的探索路径之一。
